\section{Introduction}
The level of wages and earnings of the workforce is determined by the forces of demand for and supply of labor based on the neo-classical framework. The same however is not true for developing countries, where labor market dualism and strong entry barriers across different segments of the labor market exists. The dual labor market model supposes the existence of two distinct sectors of economic activity usually classified as formal and informal sector, which extends to the nature of the employment \footnote{The distinction between the formal and informal sector on the one hand and formal and informal employment on the other has become increasingly important with growing numbers of employed persons who work in formal sector establishments also do not have access to basic benefits. The focus on informal employment is therefore to identify persons who work in precarious employment situations, irrespective of the sector into which the entity for which they work falls.} . Within formal employment, employees are subject to stable jobs with higher pay, better working conditions, permanent contract, paid parental and annual leave, paid leave for illness and accidents and social security contributions provided by the employer. Although informal employment has no employment contracts and does not ensures these jobs based benefits, it does provide more flexibility in terms of the working hour and places of work. This asymmetry affects how employees value formal and informal employment differently. 

Although those in poverty are more likely to participate in informal employment, their wage rates and total income from informal work are lower than formal works, which furthermore reinforces their marginalized position (Williams, 2014). The productivity level of informal employment is lower than the formal employment due to the differences in technology \textcite(Monsted2020) (Monsted, 2020), lack of transparency of their own accounts, long established territory based on well-entrenched territory control and rents, sub-optimal allocation of productive factors, reliance on family sources for credit, prevents acquiring modern management skills and worker training, limiting growth potential and access to world market and fragile working condition (Benjamin et al., 2014). Because of the mismatch in productivity and subsequently different wage setting mechanism, the extent of wage discrimination for the formal and informal employees is likely to vary based on personal and household, job and employer characteristics. Within formal employment, there is anti-discrimination and equal treatment laws and regulations in place, whereas these regulations and wage bargaining process are completely absent for informal employment, with competition as well as social norms favoring the dominant group (Neog \& Sahoo, 2016). Understanding of the wage discrimination with the existence of distinct employment opportunities and differing characteristics via distinct policy recognition is crucial to understand the labor market dynamics and welfare distribution within the economy.

This paper focuses on the individuals of 15-60 years who are in full time wage employment defined as individuals as who work in private financial business firm, private non-financial institutions, non-profit institutions and other institutions for 40 hours or more during the week with cash as the mode of the payment \footnote{Those who work in government, state-owned enterprises and international organizations/foreign embassy are excluded. Similarly, formal sector and informal sector both comprises of non-agriculture data and those whose economic activities are in the agricultural industries are reported on separately within job sectors.} . It aims to examine whether significant gaps in earnings exists among both formal and informal wage earners. There are several reasons motivating for the choice of full time wage employees. First, it is important to analyze employees separately from other groups, such as self-employed workers, because certain determinants of differences in earnings such as employer characteristics and job characteristics apply only to full time wage employees. Second, the definition of informality is not universal and the definition varies across employment statuses. In my study, I have used the conceptual framework of informal employment as defined by the 15th International Conference of Labour Statisticians as the Nepal Labor Force Survey III follows the same definition \footnote{15th International Conference of Labour Statisticians (ILO, 2018) defines Informal Employment as “the total number of informal jobs, whether carried out in formal sector enterprises, informal sector enterprises, or households, or as the total number of persons engaged in informal jobs during a given reference period”. In the Nepal Labour Force Survey III, those in informal employment are identified and the residual is those in formal employment.} . Third, focusing on employees makes it easier to compare and identify the gaps in literature which is concentrated mostly on gender wage gap and lacks the study based on wage gaps on informality. Fourth, excluding agriculture sector, 41 percent of total employment is based on informal sector, which compromises the single largest most sector of employment in Nepal. To complement the analysis. I provide the analysis of formal and informal employment based on 

In this paper, I use OLS regressions to assess the roles of individual and household, firm and employer characteristics in shaping hourly wages of formal and informal employees. To recover the gaps in potential wage offers, I control for these characteristics of formal and informal employment. Moreover, I use Oaxaca-Blinder Decomposition to determine whether employees in informal employment faces a wage gap in comparison to formal employment. A participation decision of workers to participate in any sector (endogenous selection) brings selectivity bias to occur and affects the formal–informal wage gap. Failing to deal with non-random selection into formal and informal jobs is a major concern as it would lead to misleading estimates of wage gaps. In order to solve the problem of sample selection, I introduce a correction term obtained from a probit model into the wage equation. Using the available data, I expect the proportion of the informal sector to all workers to influence the sectoral choice of workers without having a direct impact on wages. The Oaxaca-Blinder methodology breaks down variations in average outcomes between groups into differences in average characteristics and rewards to those characteristics. Furthermore, I carry out quantile regression to capture workers heterogeneity in returns to observed characteristics which allows us to estimate wage gap at different quantiles of wage distribution. Finally, I use Machado and Mata (MM) decomposition with a bootstrap approach, extending the traditional Oaxaca-Blinder wage decomposition on the mean across the entire wage distribution combining quantile regression and bootstrap approach to explain the significance of covariates across all distribution. Since the unconditional quantile is not the same as the integral of the conditional quantiles, authors provide a simulation-based estimator where the counterfactual distribution is constructed from the generation of a random sample.

Using the data from the nationally representative labor force surveys conducted in 2017 referred as Nepal Labor Force Survey- third round (NLFS -III), I find that those in formal employment have greater average wages than those in informal employment. 

The remainder of the paper is organized as follows. We start by discussing the impact of informality on employment and wage inequality while reviewing the related literature. Section 3 sets up the empirical model. In Section 4, we describe the data and provide descriptive statistics on inequality among wage-earners in formal and informal employment of Nepal. We then discuss the results, looking at the selection into Oaxaca-Blinder decomposition, estimation of quantile regression and MM wage decomposition based on quantile regression for formal and informal employment of Nepal. In the section 5, we conclude with possible extensions. 


Understanding the outcome that females face when they participate in the labor market is key in designing gender-equitable labor policy, especially in the global south, where females still face myriad of challenges when they want to join or are in the labor market. Globally, there has been a convergence in female participation and wage rates \parencite{WEF2020, OECD2023}. However the pace of convergence has been stagnating in advanced economies \parencite{Blau2017}. This stagnation begs the question of whether the low-hanging fruits have largely been exhausted in developed economies, leaving only politically sensitive or economically costly gender parity policy measures. If developed economies, with their greater financial strength and better institution qualities, struggle to sustain progress, it is imperative for developing countries to book-keep the factors driving the wage gap and critically examine which aspects to be addressed for avoiding the similar fate. Our contribution to this extensive literature is twofold. The first contribution of this paper is to decompose the selection-adjusted wage distribution over two decades in a developing country, Nepal. The second is to understand important channels of structural bias, differential time allocation in home production and intra-household gender education gap, against females.\par


Generally, women, either out of societal expectations or personal choice, allocate a larger share of time in home production, which increases women's reservation wage and makes them less likely to participate in the job market. As a result, female participation in the job market invariably suffers from the sample selection issue which was first identified and addressed by \textcite{Heckman1974} and \textcite{Gronau1974}. Thus, women who participate in the job market may not represent the overall female working age population as only females with a certain set of characteristics may join the job market. Several approaches have been developed in the literature  to address the selection bias \parencite{Buchinsky1998, Blau2006, Blundell2007, Mulligan2008, Huber2015, Bar2015, Maasoumi2017, Arellano2017}. We opt for a quantile-coupula approach of \textcite{Arellano2017} and jointly estimate wage and participation equations after explicitly modeling correlation between unobservables of both equations. Methodologically, this is less restrictive compared to available alternatives and use of quantile regression helps to extract the entire wage distribution. Afterwards, we employ \textcite{Chernozhukov2013} to decompose the net difference between male and female wage distributions into composition and structure effects. Prior is the wage gap resulting from the difference in individual characteristics of males and females, and latter is the difference due to varying returns of those characteristics.\par   

The three rounds of Nepal Labor Force Surveys (1998-2018) cover some of the major events that reshaped the Nepalese labor market. The decade-long armed conflict, starting in 1996 and claiming 14,242 lives \parencite{Joshi2015}, stagnated economic growth of the country. As a result, during the conflict and the post-conflict transition, the national economy struggled to accommodate the growing youth population, which fueled the rise of international labor migration \parencite{libois2016households}. The out migration rate surged post-conflict, peaking in 2013/14 and slightly declining thereafter \parencite{MoLE/GoN2013, MFES/GoN2020}, and continues to provide a sizable remittance inflow up to a quarter of gross domestic product. Immediately after the conflict, interim constitution of 2007 introduced a reservation system in public institutions for women and marginalized segment of the society\parencite{Subedi2022, Mainali2017}. In the mean time, economic structure transformed from subsistence agriculture to service sector, without developing a significant industrial base \parencite{Sapkota2013}. This change led to proliferation of service sector jobs in the newly liberalized parts of economy, mainly in education, health, and finance. Thus in the two decades of data coverage, the initial one exhibits armed conflict and minimal job market changes, whereas the second entails peak out migration, rapid growth of service sector, and the implementation of reservation system.\par

We find notable trends of gender wage convergence between high-earners, while a widening or stagnated gap among low-earners. The ``sticky floor'', rather than ``glass ceiling'' phenomenon seems to be a more apt characterization of this change. The female labor force participation declined from 30.2\% during agriculture led job market of 1998 to 16.9\% in 2018, when service sector dominated available jobs. This change in the available jobs' nature along with affirmative action policies, mandating 33\% women participation, selectively benefited educated women at the upper end of the wage distribution. The educational advantage in service sector led to a significant influx of educated women into the labor force over the course of two decades. As a result,  contribution of composition effect in the wage gap grossly vanished by 2018 across the entire wage distribution. Overall in these two decades, wage disparities remarkably shifted towards being mostly structurally driven rather than compositional one.\par 

The qualitative change in the nature of gap, i.e., very meager compositional gap but almost all structural gap, imply improving human capital alone won't budge the gap. The culprit, structural effect, comes from two sources: first, differing returns to observed characteristics and second, unobserved characteristics in the wage equation. The second part, unobserved characteristics can be any of myriad of variables that have been studied in the literature like  personal preferences\parencite{Wiswall2017, LeBarbanchon2020}, household dynamics \parencite{Bertrand2015, Goldin2017}, job characteristics \parencite{Card2015}, and societal structures \parencite{Givord2015, Lippmann2020, Goldin2006, Becker2019} over others that effect female job market participation and outcome. To understand the increasing trend of structural effects especially in urban areas, we look at the household level dynamics and check how they suppress women from participating in gainful employment. The stylized household decision-making model of \textcite{Cortes2020} makes some interesting predictions. If wives have or are presumed to have advantage in household tasks, they are less likely to participate in job market vis-a-vis their husbands. On top of it, if husbands have larger advantage in job market, this advantage skews wives more so towards household chores. \par

We test this prediction in national census (2011) by looking into the relation between differential earning potential and employment status. We find spousal education gap, the proxy for difference in earning potential, hinders female job market participation and promotes sorting of female into home production especially after marriage. The negative effect of spousal education gap increasingly overshadows gains from years of schooling when single women first marries and later becomes spouse of household head. Interestingly, the spousal education gap does not hinder female participation in own account work, but only employment in labor market.\par 

Further, we find this observed gendered sorting to be consistent with the time allocation in home production. Females allocate a similar amount of time doing household chores in three data sets:  a living standard and two labor force  surveys covering 2008 to 2018. In this decadal time frame, women's time declined by 40 minutes from previous 142 minutes, which when examined along side the employment status looks hollow. Men, regardless of employment status, contribute very little. But, females consistently work almost at the same level. The observed decline does not seem to originate from substitution across gender, but could have been through widespread adoption of home appliances and infrastructure development, like increased access to piped drinking water, that is beyond the scope of this paper. These results in combination confirm that household dynamics are important part of increasing structural effect when considering selection.\par

The paper is organized in the following fashion. In section 2, methodology, we review selection adjustment methods and describe estimation strategy used in this paper.  In section 3, we describe  nature and sources of data along with changes in the work force characteristics. Section 4 presents results and discussions of decomposition, selection, time use, and earning potential sequentially. In section 5, we conclude with possible extensions.\par                                                                                                                                                                                                                                                                                       

